% Claim proof file for row claim_3_24 -- "SigmaSeparationEquivalences".
% Auto-generated stub by scaffold/solving_current_row.py at row start. A
% manager agent dedicated to writing the TeX proof fills in the body below.
% The proof must:
%   - Stay close to the lecture notes' proof if one exists (always search
%     the LN first for a `\begin{proof}` block immediately following the
%     claim; copy / lightly edit when present).
%   - Otherwise use the LN paradigm and cite prior claims by their `ref`.
%   - Be self-contained enough that a mathematician can follow it without
%     re-reading the LN.
%   - Have no `sorry`, `omitted`, or "obvious" shortcuts.
%
% Compiles standalone via the `subfiles` package.

\documentclass[main]{subfiles}

\begin{document}
% ref: claim_3_24 (proof, with statement restated for readability)
% title: SigmaSeparationEquivalences
\def\rowref{claim\_3\_24}
\phantomsection\label{claim_3_24}
%
% Cross-references: see claim_<ref>_statement_<title>.tex in this folder
% for the corresponding statement.

% --- statement (restated) ---------------------------------------------------
\begin{Rem}[SigmaSeparationEquivalences]
    \begin{enumerate}
        \item By Proposition \refrow{claim_3_23} we have that $A \isPerp_G B \given C$ is equivalent to either of the following:
  \begin{enumerate}
      \item every walk from a node in $A$ to a node in $J \cup B$ is $C$-$\sigma$-blocked by $C$;
      \item every path from a node in $A$ to a node in $J \cup B$ is $C$-$\sigma$-blocked by $C$.
  \end{enumerate}
    \item Proposition \refrow{claim_3_23} also shows that if $A \nisPerp_G B \given C$ holds then:
  \begin{enumerate}
      \item there exists a (shortest) $C$-$\sigma$-open path from a node in $A$ to a node in $J \cup B$;
      \item there exists a (shortest) $C$-$\sigma$-open walk from a node in $A$ to a node in $J \cup B$ such that all its colliders lie in $C$.
  \end{enumerate}
  \end{enumerate}
  In practice we usually check if every path is $C$-$\sigma$-blocked or not. This is because there are, in contrast to walks, only a finite number of paths in a (finite) graph. In proofs, though, it often is easier to make use of walks, since these can be concatenated into walks (while one cannot in general concatenate two paths and again obtain a path).
\end{Rem}

% --- proof ------------------------------------------------------------------
\begin{proof}
Both items are direct corollaries of Proposition \refrow{claim_3_23} (the LN's $\mathtt{prp\!:\!sigma\_opens}$, recalling the three-way equivalence: for $C \subseteq J \cup V$ and $w_1, w_2 \in J \cup V$, the existence of a $C$-$\sigma$-open path between $w_1$ and $w_2$, the existence of a $C$-$\sigma$-open walk between $w_1$ and $w_2$, and the existence of a $C$-$\sigma$-open walk between $w_1$ and $w_2$ all of whose colliders lie in $C$ are mutually equivalent). We treat the two items in turn.

\smallskip
\noindent\textit{A remark on ``(shortest)''.}\;
The LN's parenthetical ``(shortest)'' in clauses 2(a) and 2(b) is a strengthening of plain existence that is not part of the substantive content of the present equivalences. Given any $C$-$\sigma$-open path from a node in $A$ to a node in $J \cup B$ (resp.\ any $C$-$\sigma$-open walk from a node in $A$ to a node in $J \cup B$ with all colliders in $C$), the set
\[ \big\{ \ell \in \mathbb{N} : \text{there is such a witness of length } \ell \big\} \]
is a non-empty subset of $\mathbb{N}$ and therefore (by the well-ordering of $\mathbb{N}$, equivalently $\mathtt{Nat.find}$ on the length-of-witness predicate) admits a least element. The corresponding witness is a shortest one. We therefore prove plain existence below; the strengthening to ``shortest'' follows by this one-line well-foundedness argument and is deferred to downstream consumers that actually need length-minimality.

\medskip
\noindent\textbf{Item 1.}\;
We prove that the following three statements are mutually equivalent:
\begin{enumerate}\setlength{\itemsep}{2pt}
  \item[(0)] $A \isPerp_G B \given C$;
  \item[(1)] every walk from a node in $A$ to a node in $J \cup B$ is $C$-$\sigma$-blocked by $C$;
  \item[(2)] every path from a node in $A$ to a node in $J \cup B$ is $C$-$\sigma$-blocked by $C$.
\end{enumerate}
It suffices to establish $(0) \iff (1)$ and $(1) \iff (2)$; transitivity of $\iff$ then closes all three pairs.

\smallskip
\noindent\emph{$(0) \iff (1)$.}\;
This is the defining clause of $A \isPerp_G B \given C$ (see \refrow{def_3_18}, clause~(1)): by definition, $A \isPerp_G B \given C$ \emph{means} ``every walk from a node in $A$ to a node in $J \cup B$ is $C$-$\sigma$-blocked by $C$''. The biconditional is therefore an unfolding of the defining equation, not a content-bearing step; we record it as a separate sub-equivalence only because the LN's clause 1(a) names this universal explicitly alongside the predicate.

\smallskip
\noindent\emph{$(1) \implies (2)$.}\;
A path is in particular a walk, so the universal over walks restricts to the universal over paths. Concretely, suppose $(1)$ holds and let $\pi$ be any path from a node $v \in A$ to a node $w \in J \cup B$. Then $\pi$ is in particular a walk from $v \in A$ to $w \in J \cup B$, hence by $(1)$ it is $C$-$\sigma$-blocked by $C$. Since the path $\pi$ was arbitrary, $(2)$ holds.

\smallskip
\noindent\emph{$(2) \implies (1)$.}\;
By contrapositive. Suppose $(1)$ fails, i.e.\ there exist $v \in A$, $w \in J \cup B$, and a walk $\pi$ from $v$ to $w$ that is \emph{not} $C$-$\sigma$-blocked by $C$ -- equivalently, $\pi$ is $C$-$\sigma$-open. Apply \refrow{claim_3_23} (the implication from clause~(2) to clause~(1) of $\mathtt{prp\!:\!sigma\_opens}$, instantiated at $w_1 := v$ and $w_2 := w$, noting $v \in A \subseteq J \cup V$ and $w \in J \cup B \subseteq J \cup V$): from the existence of a $C$-$\sigma$-open walk between $v$ and $w$ we obtain the existence of a $C$-$\sigma$-open \emph{path} between $v$ and $w$. Pick such a path $\pi'$; then $\pi'$ is a path from $v \in A$ to $w \in J \cup B$ that is not $C$-$\sigma$-blocked by $C$, contradicting $(2)$. Hence $(2) \implies (1)$, and this is the only step of item 1 that uses \refrow{claim_3_23}.

This closes item 1.

\medskip
\noindent\textbf{Item 2.}\;
We prove that the following three statements are mutually equivalent:
\begin{enumerate}\setlength{\itemsep}{2pt}
  \item[(0)] $A \nisPerp_G B \given C$;
  \item[(1)] there exists a $C$-$\sigma$-open path from a node in $A$ to a node in $J \cup B$;
  \item[(2)] there exists a $C$-$\sigma$-open walk from a node in $A$ to a node in $J \cup B$ such that all its colliders lie in $C$.
\end{enumerate}
Note the deliberate asymmetry between $(1)$ and $(2)$: clause $(1)$ asks only for a $C$-$\sigma$-open path (no collider-control conjunct), while clause $(2)$ adds the conjunct ``all colliders of the witness walk lie in $C$''. This mirrors clauses~(1) and~(3) of $\mathtt{prp\!:\!sigma\_opens}$ exactly and is what \refrow{claim_3_23} delivers between them; we preserve it. Again it suffices to establish $(0) \iff (1)$ and $(1) \iff (2)$.

\smallskip
\noindent\emph{$(0) \implies (1)$.}\;
By \refrow{def_3_18}, clause~(2), $A \nisPerp_G B \given C$ is the LN's shorthand for the negation $\neg \lp A \isPerp_G B \given C \rp$. Unfolding $\isPerp_G$ via clause~(1) of the same definition, this negation reads: \emph{not} every walk from a node in $A$ to a node in $J \cup B$ is $C$-$\sigma$-blocked. Pushing the negation classically through the three universal quantifiers (one each for $v$, $w$, and the walk $\pi$), there exist $v \in A$, $w \in J \cup B$, and a walk $\pi$ from $v$ to $w$ such that $\pi$ is \emph{not} $C$-$\sigma$-blocked by $C$ -- equivalently, $\pi$ is $C$-$\sigma$-open. Now apply \refrow{claim_3_23} (the implication from clause~(2) to clause~(1) of $\mathtt{prp\!:\!sigma\_opens}$, with $w_1 := v$ and $w_2 := w$): from the $C$-$\sigma$-open walk $\pi$ between $v$ and $w$ we obtain a $C$-$\sigma$-open \emph{path} $\pi'$ between $v$ and $w$. The triple $(v, w, \pi')$ witnesses $(1)$: $v \in A$, $w \in J \cup B$, and $\pi'$ is a $C$-$\sigma$-open path from $v$ to $w$.

\smallskip
\noindent\emph{$(1) \implies (0)$.}\;
Suppose $(1)$: pick $v \in A$, $w \in J \cup B$, and a $C$-$\sigma$-open path $\pi$ from $v$ to $w$. A path is in particular a walk, so $\pi$ is a $C$-$\sigma$-open walk from $v \in A$ to $w \in J \cup B$. The triple $(v, w, \pi)$ thereby witnesses the failure of the universal ``every walk from a node in $A$ to a node in $J \cup B$ is $C$-$\sigma$-blocked'' -- there is a walk between such nodes that is not $C$-$\sigma$-blocked. Hence $A \isPerp_G B \given C$ fails, i.e.\ $A \nisPerp_G B \given C$ holds (\refrow{def_3_18}, clause~(2)), which is $(0)$. This direction does not require \refrow{claim_3_23}; it is the trivial path-to-walk weakening composed with the definition of $\nisPerp_G$.

\smallskip
\noindent\emph{$(1) \iff (2)$.}\;
This is a direct application of \refrow{claim_3_23}, packaged inside the surrounding existentials $\exists v \in A, \exists w \in J \cup B$.

Concretely, fix any $v \in A \subseteq J \cup V$ and $w \in J \cup B \subseteq J \cup V$. By \refrow{claim_3_23} (the equivalence between clauses~(1) and~(3) of $\mathtt{prp\!:\!sigma\_opens}$, instantiated at $w_1 := v$ and $w_2 := w$),
\[
  \big( \exists \text{ $C$-$\sigma$-open path between $v$ and $w$ in $G$} \big)
  \;\iff\;
  \big( \exists \text{ $C$-$\sigma$-open walk between $v$ and $w$ in $G$ with all colliders in $C$} \big).
\]
This is a pointwise (in $v$ and $w$) equivalence. Quantifying both sides existentially over $v \in A$ and $w \in J \cup B$, the existential quantifiers distribute through the equivalence:
\[
  \big( \exists v \in A,\, \exists w \in J \cup B,\, \exists \text{ $C$-$\sigma$-open path from $v$ to $w$} \big)
  \iff
  \big( \exists v \in A,\, \exists w \in J \cup B,\, \exists \text{ $C$-$\sigma$-open walk from $v$ to $w$ with colliders in $C$} \big).
\]
The left-hand side is $(1)$ and the right-hand side is $(2)$. (The asymmetry flagged above is intact: the path side of $\mathtt{prp\!:\!sigma\_opens}$ carries no collider-control conjunct -- clause~(1) -- and the walk side carries the full collider-control conjunct -- clause~(3) -- and we transport both unchanged.)

This closes item 2, and the proof is complete.
\end{proof}

\end{document}
